\chapter{CƠ SỞ LÝ THUYẾT}
\label{chap:lythuyet}

Chương này trình bày nền tảng của các thành phần trong hệ thống: mạng nơ-ron tích chập và bài toán phát hiện đối
tượng, hai họ mô hình phát hiện được so sánh (YOLO và DETR), các thuật toán theo dõi đa đối tượng họ SORT, và các chỉ
số đánh giá.

\section{Mạng nơ-ron tích chập và bài toán phát hiện đối tượng}

\subsection{Mạng nơ-ron tích chập}

Mạng nơ-ron tích chập (CNN) xử lý ảnh bằng các lớp tích chập: mỗi lớp trượt một tập bộ lọc nhỏ trên bản đồ đặc trưng
đầu vào. Với đầu vào $X$ có $C$ kênh, bộ lọc kích thước $k \times k$ và bước trượt $s$, kênh thứ $c'$ của đầu ra là
\begin{equation}
  Y(i, j, c') = \sum_{c=1}^{C} \sum_{u=0}^{k-1} \sum_{v=0}^{k-1} W(u, v, c, c')\, X(i s + u,\; j s + v,\; c) + b_{c'}
  \label{eq:conv}
\end{equation}
Công thức (\ref{eq:conv}): phép tích chập hai chiều.

Với:
\begin{itemize}
  \item $W$: trọng số bộ lọc, dùng chung cho mọi vị trí trong ảnh;
  \item $b_{c'}$: độ lệch của kênh đầu ra $c'$;
  \item $s$: bước trượt; $s = 2$ làm bản đồ đặc trưng nhỏ đi một nửa mỗi chiều.
\end{itemize}

Vì trọng số dùng chung theo vị trí, số tham số không phụ thuộc kích thước ảnh và mạng nhận ra cùng một mẫu (bánh xe,
kính chắn gió) ở mọi chỗ trong ảnh. Xếp chồng nhiều lớp tích chập, chuẩn hoá theo batch và hàm kích hoạt phi tuyến (các
mô hình YOLO gần đây dùng SiLU, $\mathrm{SiLU}(x) = x \cdot \sigma(x)$) giúp vùng nhìn (receptive field) của mỗi nơ-ron
lớn dần: lớp nông bắt cạnh và góc, lớp sâu bắt hình dạng cả chiếc xe. Mạng ResNet~\cite{he2016resnet} thêm kết nối tắt
$y = F(x) + x$ để huấn luyện được mạng sâu; RT-DETRv2-S dùng ResNet-18 làm phần trích đặc trưng.

\subsection{Bài toán phát hiện đối tượng}

Mô hình phát hiện nhận một ảnh và trả về tập dự đoán $\{(b_k, c_k, s_k)\}$, trong đó $b_k = (x, y, w, h)$ là box (toạ độ
tâm, chiều rộng, chiều cao), $c_k$ là lớp và $s_k$ là điểm tin cậy. Các mô hình hiện đại thường gồm ba phần:
\begin{itemize}
  \item \textbf{Backbone}: mạng trích đặc trưng, cho ra các bản đồ đặc trưng ở nhiều độ phân giải (thường là 1/8, 1/16 và
        1/32 ảnh gốc);
  \item \textbf{Neck}: trộn đặc trưng giữa các mức để mức độ phân giải cao (dành cho vật nhỏ) cũng mang thông tin ngữ
        nghĩa của mức sâu;
  \item \textbf{Head}: dự đoán box và lớp tại từng vị trí trên các bản đồ đặc trưng.
\end{itemize}

Mức độ trùng khớp giữa hai box $A$ và $B$ được đo bằng tỉ lệ giao trên hợp (IoU):
\begin{equation}
  \mathrm{IoU}(A, B) = \frac{|A \cap B|}{|A \cup B|}
  \label{eq:iou}
\end{equation}
Công thức (\ref{eq:iou}): tỉ lệ giao trên hợp của hai box, nhận giá trị từ 0 (không giao nhau) đến 1 (trùng khít).

IoU được dùng ở ba chỗ: khi gán nhãn huấn luyện (dự đoán nào chịu trách nhiệm cho xe nào), khi chấm điểm (một dự đoán
là đúng nếu IoU với box nhãn vượt ngưỡng) và khi theo dõi (ghép box giữa hai frame).

Nhiều mô hình phát hiện, kể cả các bản YOLO trước YOLOv10, cho ra nhiều box chồng lên cùng một xe. Bước hậu xử lý
\emph{loại bỏ không cực đại} (NMS) giữ box có điểm cao nhất rồi xoá các box khác có IoU với nó vượt ngưỡng. NMS có
một tham số phải chỉnh tay và hoạt động kém ở cảnh xe đứng sát nhau: hai xe máy dừng cạnh nhau có IoU cao, NMS có thể
xoá nhầm một xe. Cả YOLO26 và họ DETR dùng trong đề tài đều không cần NMS.

\section{Họ mô hình YOLO và YOLO26}

\subsection{Ý tưởng phát hiện một giai đoạn}

YOLO~\cite{redmon2016yolo} coi phát hiện đối tượng là một bài toán hồi quy: một lần truyền xuôi qua mạng cho ra trực
tiếp toạ độ box và xác suất lớp tại mọi vị trí, thay cho cách hai giai đoạn (đề xuất vùng rồi phân loại từng vùng). Cách
làm này nhanh, phù hợp xử lý video. Các phiên bản sau giữ ý tưởng đó và cải tiến dần backbone, neck đa tỉ lệ, cách gán
nhãn và hàm mất mát.

\subsection{YOLO26}

YOLO26~\cite{jocher2026yolo26} là thế hệ YOLO mới nhất của Ultralytics~\cite{ultralytics}, có năm cỡ n, s, m, l, x. Đề
tài dùng cỡ n (2.5 triệu tham số) làm mô hình triển khai và cỡ s (10.0 triệu tham số) để so sánh. Các điểm khác biệt
chính so với thế hệ trước:
\begin{itemize}
  \item \textbf{Phát hiện đầu cuối, không cần NMS.} Theo hướng của YOLOv10~\cite{wang2024yolov10}, khi suy luận mô hình
        dùng đầu ra gán nhãn một-một: mỗi xe chỉ ứng với một dự đoán, nên đầu ra dùng được ngay, không qua NMS.
  \item \textbf{Bỏ Distribution Focal Loss (DFL).} Box được hồi quy trực tiếp, giúp mô hình gọn hơn khi xuất sang các
        định dạng triển khai như OpenVINO.
  \item \textbf{ProgLoss và STAL.} ProgLoss cân bằng dần các thành phần mất mát trong quá trình huấn luyện; STAL là cách
        gán nhãn chú ý tới vật nhỏ. Cả hai nhắm vào vật nhỏ, đúng với xe máy ở xa trong bài toán này.
  \item \textbf{Optimizer MuSGD}, kết hợp SGD với Muon, dùng cho các lần huấn luyện dài.
\end{itemize}

Nhật ký huấn luyện của đề tài ghi ba thành phần mất mát: mất mát box dựa trên IoU, mất mát phân lớp và mất mát L1 trên
toạ độ box. Mất mát phân lớp là entropy chéo nhị phân (BCE) tính riêng cho từng lớp:
\begin{equation}
  \mathcal{L}_{\mathrm{cls}} = -\frac{1}{N} \sum_{i=1}^{N} \sum_{c=1}^{K}
  \Big[ y_{ic} \log p_{ic} + (1 - y_{ic}) \log (1 - p_{ic}) \Big]
  \label{eq:bce}
\end{equation}
Công thức (\ref{eq:bce}): mất mát phân lớp entropy chéo nhị phân.

Với:
\begin{itemize}
  \item $N$: số dự đoán được tính mất mát; $K = 4$: số lớp;
  \item $y_{ic}$: nhãn đích của dự đoán $i$ với lớp $c$;
  \item $p_{ic}$: xác suất mô hình dự đoán (sau hàm sigmoid).
\end{itemize}

Tổng mất mát là tổng có trọng số $\mathcal{L} = \lambda_{\mathrm{box}} \mathcal{L}_{\mathrm{box}} +
\lambda_{\mathrm{cls}} \mathcal{L}_{\mathrm{cls}} + \lambda_{L1} \mathcal{L}_{L1}$, với các hệ số lấy theo mặc định
của thư viện.

\textbf{Ưu điểm:} nhanh, nhẹ, có sẵn công cụ huấn luyện, xuất mô hình và theo dõi đối tượng trong cùng thư viện;
không cần NMS. \textbf{Nhược điểm:} ảnh được thu nhỏ về 960 điểm ảnh cạnh dài (với cấu hình của đề tài), nên xe rất
nhỏ ở xa mất chi tiết; cỡ n có ít tham số nên kém hơn cỡ lớn ở cảnh khó.

\section{Transformer và họ mô hình DETR}

\subsection{Cơ chế attention}

Transformer~\cite{vaswani2017attention} xây dựng trên cơ chế attention: mỗi phần tử của chuỗi đầu vào tổng hợp thông
tin từ mọi phần tử khác theo trọng số học được.
\begin{equation}
  \mathrm{Attention}(Q, K, V) = \mathrm{softmax}\!\left( \frac{Q K^{\top}}{\sqrt{d_k}} \right) V
  \label{eq:attention}
\end{equation}
Công thức (\ref{eq:attention}): attention tích vô hướng có chuẩn hoá.

Với:
\begin{itemize}
  \item $Q$, $K$, $V$: ma trận truy vấn (query), khoá (key) và giá trị (value), là các phép chiếu tuyến tính của đầu
        vào;
  \item $d_k$: số chiều của khoá; chia cho $\sqrt{d_k}$ để hàm softmax không bị bão hoà.
\end{itemize}

Attention nhiều đầu (multi-head) chạy song song nhiều phép attention với các phép chiếu khác nhau rồi ghép kết quả.
Khác với tích chập chỉ nhìn một vùng lân cận, attention nhìn toàn ảnh, nhưng chi phí tăng theo bình phương số phần tử.

\subsection{DETR: phát hiện như bài toán dự đoán tập hợp}

DETR~\cite{carion2020detr} dùng CNN trích đặc trưng, đưa qua bộ mã hoá và bộ giải mã Transformer với $N$ truy vấn đối
tượng (object query); mỗi truy vấn cho ra một box và một lớp, có thể là lớp ``không có đối tượng'' ($\varnothing$).
Khi huấn luyện, $N$ dự đoán được ghép một-một với các box nhãn bằng thuật toán Hungarian:
\begin{equation}
  \begin{split}
    \hat{\sigma} &= \arg\min_{\sigma} \sum_{i=1}^{N} \mathcal{L}_{\mathrm{match}}\big(y_i, \hat{y}_{\sigma(i)}\big) \\
    \mathcal{L}_{\mathrm{match}} &= -\mathbf{1}_{\{c_i \neq \varnothing\}}\, \hat{p}_{\sigma(i)}(c_i)
      + \mathbf{1}_{\{c_i \neq \varnothing\}}\, \mathcal{L}_{\mathrm{box}}\big(b_i, \hat{b}_{\sigma(i)}\big)
  \end{split}
  \label{eq:detr-match}
\end{equation}
Công thức (\ref{eq:detr-match}): ghép cặp tối ưu giữa nhãn và dự đoán trong DETR.

Với:
\begin{itemize}
  \item $\sigma$: một hoán vị gán mỗi nhãn $y_i = (c_i, b_i)$ cho một dự đoán $\hat{y}_{\sigma(i)}$;
  \item $\hat{p}_{\sigma(i)}(c_i)$: xác suất dự đoán $\sigma(i)$ gán cho lớp đúng $c_i$;
  \item $\mathcal{L}_{\mathrm{box}} = \lambda_{L1} \lVert b_i - \hat{b}_{\sigma(i)} \rVert_1 +
        \lambda_{\mathrm{giou}} \mathcal{L}_{\mathrm{giou}}$: mất mát box.
\end{itemize}

Thành phần GIoU khắc phục việc IoU bằng 0 cho mọi cặp box không giao nhau:
\begin{equation}
  \mathrm{GIoU}(A, B) = \mathrm{IoU}(A, B) - \frac{|C| - |A \cup B|}{|C|},
  \qquad \mathcal{L}_{\mathrm{giou}} = 1 - \mathrm{GIoU}
  \label{eq:giou}
\end{equation}
Công thức (\ref{eq:giou}): IoU tổng quát, với $C$ là box nhỏ nhất bao cả $A$ và $B$.

Vì ghép một-một, mỗi xe chỉ ứng với một dự đoán và DETR không cần NMS. Nhược điểm của DETR gốc là hội tụ chậm (hàng
trăm epoch trên COCO) và chậm khi suy luận.

\subsection{RT-DETR và RT-DETRv2}

RT-DETR~\cite{zhao2024rtdetr} đưa họ DETR về tốc độ thời gian thực trên GPU nhờ ba thay đổi:
\begin{itemize}
  \item \textbf{Bộ mã hoá lai.} Attention trong cùng một mức (AIFI) chỉ chạy trên bản đồ đặc trưng sâu nhất (1/32),
        nơi số phần tử ít; việc trộn đặc trưng giữa các mức (CCFM) làm bằng tích chập.
  \item \textbf{Chọn truy vấn theo IoU.} Truy vấn khởi tạo cho bộ giải mã được chọn từ các vị trí có điểm phân lớp và
        IoU dự kiến cao, thay cho truy vấn học tự do.
  \item \textbf{Attention biến dạng trong bộ giải mã.} Mỗi truy vấn chỉ lấy mẫu một số ít điểm quanh điểm tham chiếu:
\end{itemize}
\begin{equation}
  \mathrm{DeformAttn}(z_q, p_q, x) = \sum_{m=1}^{M} W_m \Big[ \sum_{k=1}^{K} A_{mqk} \cdot W'_m\, x\big(p_q +
  \Delta p_{mqk}\big) \Big]
  \label{eq:deform}
\end{equation}
Công thức (\ref{eq:deform}): attention biến dạng.

Với:
\begin{itemize}
  \item $z_q$, $p_q$: đặc trưng và điểm tham chiếu của truy vấn $q$;
  \item $M$: số đầu attention; $K$: số điểm lấy mẫu mỗi đầu;
  \item $\Delta p_{mqk}$: độ lệch vị trí học được; $A_{mqk}$: trọng số attention, tổng theo $k$ bằng 1.
\end{itemize}

Mất mát phân lớp của RT-DETR là Varifocal Loss (VFL), dùng IoU giữa box dự đoán và box nhãn làm nhãn mềm $q$:
\begin{equation}
  \mathrm{VFL}(p, q) =
  \begin{cases}
    -q \big( q \log p + (1 - q) \log (1 - p) \big), & q > 0 \\
    -\alpha\, p^{\gamma} \log (1 - p), & q = 0
  \end{cases}
  \label{eq:vfl}
\end{equation}
Công thức (\ref{eq:vfl}): Varifocal Loss, với $p$ là điểm dự đoán, $\alpha$ và $\gamma$ là hệ số giảm trọng số cho
mẫu nền.

RT-DETRv2~\cite{lv2024rtdetrv2} giữ kiến trúc trên và thêm một số cải tiến không làm chậm suy luận (bag-of-freebies):
số điểm lấy mẫu khác nhau cho từng mức đặc trưng, toán tử lấy mẫu rời rạc thuận tiện khi triển khai, tăng cường dữ
liệu động và siêu tham số điều chỉnh theo cỡ mô hình. Khi huấn luyện còn dùng thêm truy vấn nhiễu (denoising) để bộ giải
mã hội tụ nhanh hơn. Bản RT-DETRv2-S dùng backbone ResNet-18 và có 20.1 triệu tham số.

\textbf{Ưu điểm:} box khít, không cần NMS, nhìn được ngữ cảnh toàn ảnh. \textbf{Nhược điểm:} nặng hơn YOLO cùng thế
hệ; trên CPU chậm hơn nhiều lần (mục~\ref{sec:toc-do}).

\section{Theo dõi đa đối tượng}

\subsection{Theo dõi dựa trên phát hiện}

Các thuật toán theo dõi trong đề tài đều theo vòng lặp sau cho mỗi frame $t$:
\begin{enumerate}
  \item chạy mô hình phát hiện, được tập box $D_t$;
  \item dự đoán vị trí ở frame $t$ của mọi quỹ đạo (track) đang theo dõi bằng bộ lọc Kalman;
  \item ghép box $D_t$ với các track dựa trên IoU giữa box và vị trí dự đoán;
  \item cập nhật track được ghép, tạo track mới cho box chưa ghép đủ tin cậy, giữ track chưa ghép ở trạng thái ``mất''
        thêm một số frame trước khi xoá.
\end{enumerate}

\subsection{Bộ lọc Kalman}

Bộ lọc Kalman~\cite{kalman1960} ước lượng trạng thái của từng xe dưới giả thiết vận tốc không đổi. Với BoT-SORT, trạng
thái là $\mathbf{x} = [x_c, y_c, w, h, \dot{x}_c, \dot{y}_c, \dot{w}, \dot{h}]^{\top}$ (tâm, kích thước box và vận tốc
của chúng). Mỗi frame gồm hai bước:
\begin{equation}
  \hat{\mathbf{x}}_{t|t-1} = F \hat{\mathbf{x}}_{t-1}, \qquad P_{t|t-1} = F P_{t-1} F^{\top} + Q
  \label{eq:kalman-predict}
\end{equation}
Công thức (\ref{eq:kalman-predict}): bước dự đoán.
\begin{equation}
  \begin{split}
    K_t &= P_{t|t-1} H^{\top} \big( H P_{t|t-1} H^{\top} + R \big)^{-1} \\
    \hat{\mathbf{x}}_{t} &= \hat{\mathbf{x}}_{t|t-1} + K_t \big( \mathbf{z}_t - H \hat{\mathbf{x}}_{t|t-1} \big) \\
    P_t &= (I - K_t H) P_{t|t-1}
  \end{split}
  \label{eq:kalman-update}
\end{equation}
Công thức (\ref{eq:kalman-update}): bước cập nhật khi track được ghép với một box.

Với:
\begin{itemize}
  \item $F$: ma trận chuyển trạng thái (vị trí cộng vận tốc); $H$: ma trận quan sát, lấy 4 thành phần box từ trạng thái;
  \item $P$: hiệp phương sai của ước lượng; $Q$, $R$: nhiễu của mô hình chuyển động và nhiễu đo;
  \item $\mathbf{z}_t$: box được ghép ở frame $t$; $K_t$: hệ số Kalman, cân bằng giữa dự đoán và quan sát.
\end{itemize}

Khi xe bị che vài frame, track không được cập nhật nhưng bước dự đoán vẫn đẩy vị trí theo vận tốc, nên khi xe xuất hiện
lại vẫn có thể ghép đúng.

\subsection{Ghép cặp bằng thuật toán Hungarian}

Giữa $n$ track và $m$ box, ma trận chi phí có phần tử $c_{ij} = 1 - \mathrm{IoU}(\hat{b}_i, d_j)$. Thuật toán
Hungarian~\cite{kuhn1955hungarian} tìm phép gán một-một có tổng chi phí nhỏ nhất:
\begin{equation}
  \min_{x_{ij} \in \{0, 1\}} \sum_{i=1}^{n} \sum_{j=1}^{m} c_{ij} x_{ij}
  \quad \text{sao cho} \quad \sum_{j} x_{ij} \le 1,\; \sum_{i} x_{ij} \le 1
  \label{eq:hungarian}
\end{equation}
Công thức (\ref{eq:hungarian}): bài toán gán. Các cặp có chi phí vượt ngưỡng ghép bị loại dù được thuật toán chọn.

\subsection{Các thuật toán theo dõi họ SORT}

\textbf{SORT}~\cite{bewley2016sort} là dạng cơ bản: Kalman để dự đoán, Hungarian theo IoU để ghép. Track mới chỉ được
xác nhận sau vài frame liên tiếp có box (\texttt{min\_hits}) và bị xoá sau một số frame không được ghép.

\textbf{ByteTrack}~\cite{zhang2022bytetrack} giữ lại cả box điểm thấp. Box được chia hai nhóm theo điểm tin cậy: lần
ghép thứ nhất chỉ dùng box điểm cao; các track còn lại được ghép lần hai với box điểm thấp. Xe bị che một phần thường
có điểm thấp, nên lần ghép thứ hai giúp track không bị đứt. Chỉ box điểm cao chưa ghép mới tạo track mới. Vì lý do
này, ngưỡng tin cậy của mô hình phát hiện trong hệ thống được đặt ở 0.1 thay cho 0.25 thông thường: các box từ 0.1 đến
0.25 không tạo track mới nhưng được dùng ở lần ghép thứ hai.

\textbf{BoT-SORT}~\cite{aharon2022botsort} dựa trên ByteTrack và thêm ba cải tiến: trạng thái Kalman dùng trực tiếp
chiều rộng và chiều cao box; bù chuyển động camera (GMC) bằng cách ước lượng phép biến đổi giữa hai frame từ luồng
quang học thưa rồi áp vào vị trí dự đoán; và tuỳ chọn kết hợp đặc trưng ngoại hình (ReID) với IoU khi ghép.

\textbf{OC-SORT}~\cite{cao2023ocsort} xử lý trường hợp xe bị che lâu: khi track được ghép lại sau thời gian mất, thuật
toán dựng lại quỹ đạo ảo từ quan sát cuối để sửa sai số tích luỹ của Kalman, và dùng hướng chuyển động thực tế làm
thêm một thành phần chi phí khi ghép.

\section{Các chỉ số đánh giá}
\label{sec:chi-so}

\subsection{Chỉ số phát hiện: Precision, Recall, AP và mAP}

Một dự đoán là đúng (TP) nếu cùng lớp và có IoU với một box nhãn chưa được ghép vượt ngưỡng; các dự đoán được xét lần
lượt theo điểm tin cậy giảm dần. Dự đoán không ghép được là phát hiện thừa (FP), box nhãn không được ghép là bỏ sót (FN).
\begin{equation}
  \mathrm{Precision} = \frac{TP}{TP + FP}, \qquad \mathrm{Recall} = \frac{TP}{TP + FN}
  \label{eq:pr}
\end{equation}
Công thức (\ref{eq:pr}): độ chính xác và độ bao phủ.

Quét ngưỡng tin cậy từ cao xuống thấp cho đường cong Precision–Recall. Độ chính xác trung bình (AP) theo quy ước
COCO~\cite{lin2014coco} là trung bình độ chính xác nội suy tại 101 mức recall:
\begin{equation}
  \mathrm{AP} = \frac{1}{101} \sum_{r \in \{0, 0.01, \ldots, 1\}} p_{\mathrm{interp}}(r), \qquad
  p_{\mathrm{interp}}(r) = \max_{\tilde{r} \ge r} p(\tilde{r})
  \label{eq:ap}
\end{equation}
Công thức (\ref{eq:ap}): AP nội suy 101 điểm của COCO.

mAP là trung bình AP trên các lớp. mAP50 tính ở ngưỡng IoU 0.5; mAP50-95 lấy trung bình thêm trên 10 ngưỡng IoU 0.50,
0.55, \ldots, 0.95, nên đòi hỏi box khít hơn. Trong đề tài, \textbf{mAP50 chính} là trung bình AP50 của xe máy và ô tô:
các tập test Việt Nam có rất ít xe buýt và xe tải (tập test Hà Nội chỉ có 5 xe buýt và 2 xe tải), AP của hai lớp này
nhảy giữa 0 và 1 theo từng box nên không đưa vào chỉ số chính.

Một số vùng trong ảnh có xe nhưng không thể gán nhãn rõ ràng (xe bị cây che gần hết, cụm xe quá xa). Theo quy ước vùng
bỏ qua của KITTI~\cite{geiger2012kitti} (\texttt{DontCare}) và COCO (\texttt{iscrowd}), các vùng này được gán lớp
\texttt{ignore}: dự đoán rơi vào vùng ignore không bị tính là thừa, và xe trong vùng không bị tính là bỏ sót.

\subsection{MOTA}

MOTA~\cite{bernardin2008clear} gộp ba loại lỗi trên toàn bộ video:
\begin{equation}
  \mathrm{MOTA} = 1 - \frac{\sum_{t} \big( \mathrm{FN}_t + \mathrm{FP}_t + \mathrm{IDSW}_t \big)}{\sum_{t} \mathrm{GT}_t}
  \label{eq:mota}
\end{equation}
Công thức (\ref{eq:mota}): độ chính xác theo dõi đa đối tượng.

Với:
\begin{itemize}
  \item $\mathrm{FN}_t$, $\mathrm{FP}_t$: số box bỏ sót và thừa ở frame $t$;
  \item $\mathrm{IDSW}_t$: số lần một xe bị đổi ID so với frame trước;
  \item $\mathrm{GT}_t$: số box nhãn ở frame $t$.
\end{itemize}

MOTA bị chi phối bởi FN và FP, tức chủ yếu đo chất lượng phát hiện; số lần đổi ID thường nhỏ so với hai thành phần kia.

\subsection{IDF1}

IDF1~\cite{ristani2016idf1} đo mức giữ đúng danh tính. Toàn bộ quỹ đạo nhãn và quỹ đạo dự đoán được ghép một-một sao
cho số frame trùng khớp lớn nhất; IDTP là số box thuộc cặp quỹ đạo được ghép và trùng khớp, IDFP và IDFN là phần còn lại.
\begin{equation}
  \begin{gathered}
    \mathrm{IDF1} = \frac{2\,\mathrm{IDTP}}{2\,\mathrm{IDTP} + \mathrm{IDFP} + \mathrm{IDFN}} \\
    \mathrm{IDP} = \frac{\mathrm{IDTP}}{\mathrm{IDTP} + \mathrm{IDFP}}, \qquad
    \mathrm{IDR} = \frac{\mathrm{IDTP}}{\mathrm{IDTP} + \mathrm{IDFN}}
  \end{gathered}
  \label{eq:idf1}
\end{equation}
Công thức (\ref{eq:idf1}): IDF1, độ chính xác danh tính IDP và độ bao phủ danh tính IDR.

IDR chỉ phụ thuộc vào quỹ đạo nhãn được theo đúng bao nhiêu, không bị phạt khi hệ thống phát hiện thêm xe không có
nhãn. Tính chất này được dùng khi đánh giá trên tập VnTraffic có nhãn thiếu (mục~\ref{sec:theo-doi}).

\subsection{HOTA}

HOTA~\cite{luiten2021hota} tách đánh giá thành hai phần, phát hiện (DetA) và liên kết (AssA), rồi lấy trung bình nhân.
Tại ngưỡng IoU $\alpha$, sau khi ghép box nhãn với box dự đoán:
\begin{equation}
  \begin{gathered}
    \mathrm{DetA}_{\alpha} = \frac{|TP|}{|TP| + |FN| + |FP|} \\
    \mathrm{AssA}_{\alpha} = \frac{1}{|TP|} \sum_{c \in TP} \frac{|TPA(c)|}{|TPA(c)| + |FNA(c)| + |FPA(c)|}
  \end{gathered}
  \label{eq:deta-assa}
\end{equation}
\begin{equation}
  \mathrm{HOTA}_{\alpha} = \sqrt{\mathrm{DetA}_{\alpha} \cdot \mathrm{AssA}_{\alpha}}, \qquad
  \mathrm{HOTA} = \frac{1}{19} \sum_{\alpha \in \{0.05, 0.10, \ldots, 0.95\}} \mathrm{HOTA}_{\alpha}
  \label{eq:hota}
\end{equation}
Công thức (\ref{eq:deta-assa}) và (\ref{eq:hota}): các thành phần của HOTA.

Với:
\begin{itemize}
  \item $TPA(c)$: tập các cặp TP có cùng ID nhãn và cùng ID dự đoán với cặp $c$, tức những frame quỹ đạo được nối đúng;
  \item $FNA(c)$, $FPA(c)$: các box của cùng xe nhãn (hoặc cùng track dự đoán) với $c$ nhưng không được ghép với nhau.
\end{itemize}

Nhờ tách hai thành phần, HOTA cho biết điểm thấp là do phát hiện kém hay do nối quỹ đạo kém. Đề tài tính MOTA, IDF1
và HOTA bằng TrackEval~\cite{trackeval}, bộ chấm được dùng cho benchmark MOTChallenge~\cite{milan2016mot16}.

\subsection{Khoảng tin cậy bootstrap}
\label{sec:bootstrap}

Các tập test của đề tài nhỏ, nên một chênh lệch mAP nhỏ có thể chỉ do mẫu ảnh. Để biết chênh lệch có đáng tin không,
đề tài lấy mẫu lại có hoàn lại các ảnh của tập test 1.000 lần, tính lại mAP trên mỗi mẫu, rồi lấy phân vị 2.5\% và
97.5\% làm khoảng tin cậy 95\%~\cite{efron1993bootstrap}. Cần khoảng 1.000 lần trở lên thì hai phân vị này mới ổn định;
với vài chục lần, chúng gần như chỉ là giá trị nhỏ nhất và lớn nhất của các mẫu. Việc ghép dự đoán với nhãn được làm
một lần cho từng ảnh, mỗi lần lấy mẫu chỉ gộp lại kết quả đã ghép của các ảnh được chọn rồi tính AP theo đúng công thức
(\ref{eq:ap}), nên 1.000 lần vẫn nhanh. Hai mô hình có khoảng tin cậy chồng nhau được coi là chưa phân biệt được.

Lấy mẫu lại từng ảnh là coi các ảnh độc lập với nhau. Với frame cắt từ cùng một clip thì không đúng: hai frame cách nhau
vài giây chứa phần lớn là cùng những chiếc xe, nên tập 30 frame chứa ít thông tin hơn 30 ảnh độc lập và khoảng tin cậy
tính theo ảnh hẹp hơn thực tế. Cách chuẩn cho dữ liệu có tương quan là lấy mẫu lại theo khối~\cite{kunsch1989block}
hoặc theo cụm~\cite{field2007cluster}: chia ảnh thành các nhóm gần như độc lập với nhau rồi lấy mẫu lại cả nhóm. Đề tài
dùng thêm cách này cho mọi tập frame cắt từ video, với nhóm là các khối 10 giây liên tiếp trong cùng một clip. Độ dài khối
được chọn từ nhãn tracking của clip Phạm Văn Đồng: 90\% số xe ở trong khung hình dưới 8.2 giây (trung vị 1.3 giây),
nên hai khối cách nhau hơn 10 giây hầu như không chung chiếc xe nào. Với UA-DETRAC, mỗi video là một nhóm. Tập có dưới 5
nhóm thì không tính khoảng tin cậy theo khối, vì khi đó số cách lấy mẫu lại khác nhau quá ít.

Cả hai cách đều chỉ đo dao động do mẫu ảnh test, không đo dao động giữa các lần huấn luyện.

\section{Kết luận chương}

Chương này đã trình bày cơ sở của hai họ mô hình phát hiện được dùng trong đề tài: YOLO26 thuộc họ CNN, phát hiện đầu
cuối không cần NMS, và RT-DETRv2 thuộc họ Transformer, ghép cặp Hungarian khi huấn luyện. Phần theo dõi gồm bộ lọc
Kalman, thuật toán Hungarian và các biến thể SORT, ByteTrack, BoT-SORT, OC-SORT. Các chỉ số mAP, MOTA, IDF1 và HOTA đo
lần lượt chất lượng phát hiện, độ chính xác theo dõi tổng thể, khả năng giữ danh tính và sự cân bằng giữa phát hiện
với liên kết. Chương~\ref{chap:xaydung} mô tả cách các thành phần này được ghép thành hệ thống và huấn luyện trên dữ
liệu Việt Nam.
